% SPDX-License-Identifier: CC-BY-4.0
% Copyright (c) 2026 Martin Schr\"oder <info@swedishembedded.com>

\section{Related work}
\label{sec:related}

\paragraph{Spectral operators on functions.} Fourier neural operators
\citep{li2020fno} parameterize an integral operator by complex weights on a
truncated set of Fourier modes, with pointwise nonlinearities applied in the
physical domain between transforms. Transform Once \citep{poli2022transformonce}
argues that repeated forward and inverse transforms are wasted work and keeps
the computation in the frequency domain across layers. Our mixer borrows the
FNO layer shape but applies it to the feature axis of each token, and it
keeps the residual stream in the frequency domain throughout, as Transform
Once recommends, transforming only around the nonlinearity.

\paragraph{Fourier token mixing.} FNet \citep{leethorp2021fnet} replaces
self-attention in an encoder with a parameter-free two-dimensional Fourier
transform over the sequence and hidden dimensions, keeping its real part, and
leaves the dense feed-forward networks in place; adaptive Fourier neural
operators \citep{guibas2021afno} mix tokens with learned block-diagonal
spectral weights. Hyena \citep{poli2023hyena} builds causal long convolutions
evaluated with FFTs as an attention replacement in language models. In all
three the Fourier transform is the token mixer and has no learned per-mode
structure on the hidden axis. SpectralLM keeps causal attention as the token
mixer and puts learned spectral structure on the feature axis, where it
replaces the feed-forward network instead.

\paragraph{Structured and spectral parameterizations of weights.} Monarch
matrices \citep{dao2022monarch} replace dense layers by products of two
block-diagonal matrices, a hardware-efficient family that can represent many
common transforms. FourierFT \citep{gao2024fourierft} learns a few spectral
coefficients of a weight update instead of a low-rank update. These compress
weights while the activations stay dense; SpectralLM changes the activations'
representation.

\paragraph{Spectral factorization of the embedding.} The closest prior work is
the L-Transformer of \citet{elichi2026structured}, which reshapes token
representations into third-order tensor slices, runs attention and
feed-forward operations in a transform domain, and proves that the result is
equivalent to $p$ parallel transformers on embeddings of $1/p$ the width,
which is where its parameter reduction comes from. That equivalence is the
tensor counterpart of \cref{prop:absorption}: a transform-domain model whose
operators are dense within each slice is a set of narrower conventional
models. It is instantiated with a discrete cosine transform and evaluated on
IMDB and AG~News text classification with an encoder. SpectralLM differs in
keeping one residual stream, in coupling all modes through a signal-domain
product, and in being an autoregressive language model.

\paragraph{Fourier features.} Random Fourier feature maps help coordinate
networks fit high-frequency functions of low-dimensional inputs
\citep{tancik2020fourierfeatures}. That is a statement about inputs that are
coordinates; the coordinates of a language model's hidden state carry no
geometry, which is why \cref{sec:absorption} matters here.

\paragraph{Small language models.} TinyStories \citep{eldan2023tinystories}
showed that models of a few million parameters learn fluent English on a
corpus of simple stories, which makes it a corpus on which architecture
differences at small scale are measurable with modest compute.
